% ID: SEL-D0007-BD-E08
% ESTADO: cotejada. Segunda pasada visual completa y borrador compilado revisados.
% Fuente: Selectividad/matematicas_aplicadas/ORD_TITULAR_MATEMÁTICAS APLICADAS CCSS-JUNIO-IMPRESO.pdf
% SHA-256: c42a605501a277fe7194b964ad92ad66047be9a62dc254fa2219d756d3d71317
% Página: 2; bloque D; ejercicio original 8.

Se desea estimar la proporción de individuos mayores de edad de una localidad
que están en contra de la construcción de una central nuclear en su término
municipal. Para ello, se pregunta a 100 individuos mayores de edad de esa
localidad, elegidos de forma aleatoria, resultando que 45 de ellos rechazan
la construcción de la central.
\begin{enumerate}[label=\alph*)]
  \item \textbf{(1.5 puntos)} Calcule un intervalo de confianza al
  92\,\% para estimar la proporción real de individuos de esa localidad
  que están en contra de la construcción de la central.
  \item \textbf{(1 punto)} Suponiendo que se mantiene la misma proporción
  muestral y el mismo nivel de confianza del apartado anterior, determine
  el tamaño mínimo de la muestra que hay que tomar, para que al estimar la
  proporción de individuos de esa localidad que rechazan la construcción
  de la central, el error cometido sea inferior al 5\,\%.
\end{enumerate}
